\section{Rahmenbedingungen und Umfeld}

\subsection{Einschränkungen und Vorgaben}
Similar Slides soll in einer objektorientierten Programmiersprache entwickelt werden. Ferner soll \textit{ausschließlich} das proprietäre .pptx-Format unterstützen. Die Analyse und der Vergleich des Inhalts der Folien soll über einen Zugriff auf die XML-Struktur einer .pptx-Datei mithilfe von VisualBasic erfolgen. Das Ähnlichkeitsmaß vergleicht die Elemente einer Folie nach der folgenden Priorität: Zuerst der Titel, dann Graphiken und Bilder und zuletzt wird der tatsächliche Textinhalt beachtet.

% Erklären Sie hier, welche Vorgaben der Kunde bezüglich der Umsetzung des Projektes ausgesprochen hat. Wird beispielsweise verlangt, dass ein bestimmtest Framework oder eine Arbeitsweise vorausgesetzt werden? Müssen rechtliche Vorgaben eingehalten werden? Möglicherweise gibt es implizite Einschränkungen, wie eine Internetanbindung, Erreichbarkeit der Software oder die Verwendung von bestimmten Zahlungsdienstleistern.

% Mind. 6 Zeilen wären hier angebracht.

\subsection{Anwender}
Die Software richtet sich hauptsächlich an den Dozenten einer Veranstaltung und seine wissenschaftlichen Mitarbeiter. Aus dieser Personengruppe sollen besonders diese Angesprochen werden, die eine solide Basis an Office/PowerPoint-Erfahung vorzuweisen haben, aber wenig Zeit bei der Vorlesungs- und Übungsvorbereitung haben. 

% Beschreiben Sie hier die Mitglieder Ihrer Zielgruppe. Richtet sich die Software an Schüler, Studenten oder vielleicht Professoren? Beschreiben Sie für jedes Mitglied der Zielgruppe welche Charakteristiken diese ausmachen und wie dies die Nutzung des Endprodukts beeinflusst. Richtet sich beispielsweise Ihre Software an ältere Menschen könnte dies die Computerkompetenzen oder die Erkennung von kleinen Details auf einem Bildschirm beeinflussen.

% Mind. 6 Zeilen wären hier angebracht.

\subsection{Schnittstellen und angrenzende Systeme} 
Similar Slides muss mit dem PowerPoint-Programm interagieren, um das gezielte Einfügen einzelner Folien zu ermöglichen. Außerdem muss die Software den XML-Unterbau einer .pptx-Datei ausgelesen und analysiert können. Zuletzt muss auch noch auf ein zentrales Repository mit Lese- und Schreibvorgängen zugregriffen werden. 

% Nennen und erläutern Sie hier die verschiedenen Systeme / Libraries, die mit Ihrem Projekt interagieren müssen. Beipsielsweise verschiedene Drittanbieter, die Services wie Zahlungsabwicklung bereitstellen oder APIs, über die Sie Informationen beziehen oder Befehle ausführen müssen.